\section{模型评价}
本文的优点是：预测选型与测试严格隔离；任务侧使用增量计算但不省略最终 5 万任务全量审计；储能候选显式包含 no-action；Q4 每轮从上一轮储能结果动态重建邻域；全部结果具有结构化运行证据和两次独立复算哈希。

局限性也必须明确。第一，Q2/Q4 只在预声明有限候选预算内寻找 best-found 可行解，不具有全局最优证书；Q3 只在离散 SOC 网格和七类目标策略中做候选集 Pareto 判断。第二，题面未给储能退化成本与循环寿命上限，平衡轨迹的等效循环次数较高，不能直接作为真实设备控制指令。第三，通信模型只含固定时延，未计带宽拥塞、迁移数据量和传输能耗。第四，没有独立 Oracle 或真实部署数据，因此可信等级最高为 scenario-feasible。
